% Generated draft: requires mathematical and source review.
\documentclass[11pt]{article}
\usepackage[margin=25mm]{geometry}
\usepackage{fontspec}
\setmainfont{Latin Modern Roman}
\usepackage{amsmath,amssymb,mathtools,unicode-math}
\setmathfont{Latin Modern Math}
\usepackage{enumitem,needspace,xurl,hyperref}
\setlength{\parindent}{0pt}
\setlength{\parskip}{5pt}
\setlength{\emergencystretch}{3em}
\newcommand{\N}{\mathbb{N}}
\newcommand{\Z}{\mathbb{Z}}
\newcommand{\Q}{\mathbb{Q}}
\newcommand{\R}{\mathbb{R}}
\newcommand{\C}{\mathbb{C}}
\newcommand{\F}{\mathbb{F}}
\newcommand{\A}{\mathbb{A}}
\newcommand{\PP}{\mathbb{P}}
\newcommand{\E}{\mathbb{E}}
\newcommand{\eps}{\varepsilon}
\newcommand{\dd}{\,\mathrm{d}}
\begin{document}
\section*{UnsolvedMath ID: 701555}
\textbf{Source problem code:} catalogue entry
\subsection{Definitions and mathematical statement}
For each integer \(n\geq 2\), let the virtual pure braid group \(VP_n\) be defined by the presentation
\[
VP_n=\left\langle \lambda_{ij}\ (1\leq i,j\leq n,\ i\neq j)\ \middle|\ 
\begin{aligned}
\lambda_{ij}\lambda_{kl}&=\lambda_{kl}\lambda_{ij}
&&\text{for pairwise distinct }i,j,k,l,\\
\lambda_{ki}\lambda_{kj}\lambda_{ij}&=\lambda_{ij}\lambda_{kj}\lambda_{ki}
&&\text{for pairwise distinct }k,i,j
\end{aligned}
\right\rangle.
\]
Write \(W_n\) for the set of finite words in the letters \(\lambda_{ij}^{\pm1}\), including the empty word, and write \([w]\) for the element of \(VP_n\) represented by \(w\).

\begin{enumerate}
\item[(a)] For every integer \(n\geq4\), construct a normal form for words in \(VP_n\). Under the usual effective interpretation, this asks for an explicitly described set \(\mathcal N_n\subseteq W_n\) and a procedure assigning to every \(w\in W_n\) a word \(\operatorname{NF}_n(w)\in\mathcal N_n\), such that
\[
[\operatorname{NF}_n(w)]=[w],\qquad
\operatorname{NF}_n(u)=\operatorname{NF}_n(v)
\ \Longleftrightarrow\ [u]=[v].
\]
The source does not specify a preferred shape of normal form or a computational complexity requirement; the effective interpretation above is an explicit clarification requiring review.

\item[(b)] Determine, for each integer \(n\geq4\), whether \(VP_n\) is linear. For a field \(K\), a group \(G\) is linear over \(K\) if there are a positive integer \(d\) and an injective homomorphism \(G\hookrightarrow\operatorname{GL}_d(K)\). Under the convention that an unspecified field is existentially quantified, the question is whether
\[
\forall n\in\Z,\ n\geq4:\quad
\exists\text{ a field }K_n\ \exists d_n\in\Z_{\geq1}\ 
\exists\text{ an injective homomorphism }
VP_n\longrightarrow\operatorname{GL}_{d_n}(K_n).
\]
The source does not specify a field or characteristic. It also states that \(VP_3\) is known to be linear; this is retained as a supplied background assertion.
\end{enumerate}
\subsection{Short English statement}
For every \(n\geq4\), construct a normal form for words in the presented virtual pure braid group \(VP_n\), and determine whether \(VP_n\) admits a faithful finite-dimensional representation over a field. The source states that \(VP_3\) is linear. The type of normal form and the representation field are not specified.
\subsection{Sources}
\begin{itemize}
\item Supplied local UnsolvedMath registry record, permanent ID \(701555\), exact source problem code \texttt{catalogue entry}, title \emph{Kourovka Notebook Problem 19.7}; upstream revision \texttt{d371ee917ac9647fbb06cd6fc529deec6f3dda82}. The record is labelled \texttt{open}; this is supplied registry status, not an independently verified literature assessment.
\item \emph{The Kourovka Notebook (Unsolved Problems in Group Theory)}, original problem 19.7, proposed by V. G. Bardakov in 2018. The supplied background identifies the 19th issue (2018) and the 21st issue edition (2026), arXiv:1401.0300v48. The collection description names E. I. Khukhro and V. D. Mazurov as editors.
\item The supplied registry record was created and updated on \texttt{2026-10-16T14:00:00Z}. These are registry timestamps, not source-verification dates. No separate reference is supplied for the assertion that \(VP_3\) is linear, and no source was independently consulted for this draft.
\end{itemize}
\end{document}
